\section{Audit, corrections, and integration status}
\label{sec:audit}

\subsection{What was actually checked}

The source anchor is the repository revision \pin. All eleven archives
in the supplied incoming directory at that revision were unpacked and
indexed. The review concentrated on their theorem and question ledgers
and on the passages used by the present arguments. The manuscript
review concentrated on Gaussian depth reduction, the $S_6$ obstruction,
the revised $S_8$ candidate, and the harmonic/Gamma interfaces.
This targeted review is not a certification of all manuscript volumes
or every theorem in the incoming reports.

Three source questions have precise new outcomes:

\begin{enumerate}
\item \textbf{Ordered boundary report, research question 2.}
The positive constants in \src{sections/08_gamma_basis.tex} are
evaluated by Theorem~\ref{gn:thm:K}. Its proof uses the source's
established top-row identities and a classical integral with an
explicitly checked normalization. The matrix $T_m$ and all column
orders are stated before use. The later-index determinant in
\eqref{gn:eq:F-explicit} is a different object.

\item \textbf{Balanced Stieltjes report, research question 4.}
The homogeneous colored balancing, its natural Lerch coordinates,
its ordinary Mellin moments, and its finite root-of-unity filter
are supplied in Sections~\ref{sec:colored}--\ref{sec:color-limit}.
Uniqueness is asserted for coefficients independent of the shift
and for a decay condition across an open interval of shifts.
This is the same kind of hypothesis needed in the predecessor's
uncolored uniqueness theorem.

\item \textbf{Lattice-contact report, end of the spectral-face section
and research question Q4.}
Theorem~\ref{tc:thm:Gamma} treats exactly two distinct mass values at
every pair of positive multiplicities. It supplies a finite
dictionary per spectral term. A problem involving two unequal
displaced masses above a third base value has three clusters and
remains outside this theorem.
\end{enumerate}

The companion source manifest records the archive names and their
SHA-256 hashes. The integration notes list exact predecessor paths
and theorem labels, so a later integration can reconcile the new
results with a changed repository head.

\subsection{Corrections and qualifications}

No false mathematical theorem was identified in the particular
incoming passages used for these three developments. It would be
misleading to manufacture an erratum from an open question that
the source already labels correctly. The following distinctions
are nevertheless essential when integrating the continuation.

\paragraph{The colored generator changes.}
For a nontrivial color, $\Phi(\omega,s,a)$ has no pole at $s=1$.
The appropriate balanced kernel is \eqref{col:kernel}. Copying
the extra uncolored term $\log^{q+1}x/(q+1)$ would spoil its
large-$x$ cancellation. This is a correction to that potential
extension, not a claim that the incoming uncolored definition
is erroneous.

\paragraph{Fixed-shift uniqueness can fail.}
The open-interval clause in Theorem~\ref{col:unique} cannot simply
be deleted on the full color domain. The exceptional-color
example following that theorem exhibits simultaneous vanishing
of consecutive tail coordinates. At such a point, a logarithm
times a lower balanced kernel can preserve the required decay.

\paragraph{Constant terms depend on the chosen expression.}
The raw $\lambda$-constant term of
$M_{m,q}(\lambda,a)$ differs from that of
$e^{-a\lambda}M_{m,q}(\lambda,a)$ by $L_{m,q}a^m/m$.
Equations \eqref{col:renorm} and \eqref{col:raw-CT} retain that
term explicitly. This avoids an incorrect identification of a
color limit with an uncolored ordinary moment.

\paragraph{Positivity does not settle every later determinant.}
The positive density \eqref{gn:eq:density} evaluates the
normalizing determinant. The source's later rows involve
$\Im(x+i/2)^{2r}$ with oscillatory signs. One cannot infer all
later signs by applying positivity of a Vandermonde product to
a different matrix. In particular, no all-pairs nonvanishing
claim follows from the closed formula for $K_m$ alone.

\paragraph{Finite normal moments do not give a finite-dimensional
normal distribution.}
The statistic in Theorem~\ref{gn:thm:gaussian} is bounded at fixed
$n,\sigma$. We prove equality of moments only through degree
$n+1$, and we do not assert that this range is maximal.
The distributional limit in Corollary~\ref{gn:cor:CLT} has a
separate moment-method proof.

\paragraph{Branch and color conventions remain explicit.}
The two-cluster $\Li_{\boldsymbol a}(z)$ assigns its argument only
to the largest summation index. Colored depth-two notation uses
two independent color factors and is defined separately.
When $C<B$, the cluster formula uses the continuation on
$\C\setminus[1,\infty)$; substituting a divergent power series
at a negative argument of modulus greater than one is not valid.
Likewise, the ordinary $\omega$ series for the colored balanced
kernel is used only inside its disk of convergence before
continuation to the full slit domain.

\paragraph{A small editorial correction.}
At the pinned revision, the manuscript README repeats consecutively
the paragraph describing the complete $1{,}032$-input first-letter
$X$ family, its rank $1{,}008$, and its kernel dimension $24$.
The optional patch in the integration directory removes one
copy and leaves the mathematical statements unchanged.

\subsection{Evidence and its limits}

Analytic proofs establish the identities. The accompanying exact
and numerical calculations provide independent checks on the
algebra, normalization, and selected parameter values.

The Gamma normalization suite evaluates the original coefficient
matrices at $m=1,2,3$, the parameterized determinants at unequal
real and complex shifts, and scalar digamma/polygamma insertions.
It contains fourteen checks at $42$ decimal digits; the largest
recorded relative residual is approximately $7.1\times10^{-43}$.
The finite-Gaussian suite constructs exponential jets of ordinary
one-dimensional moment integrals and takes their determinant by
finite polynomial arithmetic. It does not use the paired Gamma
product on the left side. Its fifteen checks include the fourth
moment at $n=3$ and non-quarter shifts; the largest recorded
absolute residual is approximately $4.9\times10^{-41}$.

The cluster suite contains $448$ exact rational assertions,
comparing two different coefficient constructions and checking
cluster exchange. Four independent numerical comparisons at
$90$ digits include a displacement outside the predecessor's
local collision-series disk. The colored suite supplies $41$ exact assertions and eleven
numerical checks at $50$ digits, covering finite tail and zeta
coefficients, a simultaneous-tail-zero example, ordinary moments,
root-of-unity specializations, and an anchored primitive. Its
largest numerical residual is approximately $5.4\times10^{-37}$.
Three further $55$-digit checks evaluate a negative cluster argument
$q=-4$, its Landen transform, and the oriented anchored primitive.
The separate depth-two certificate uses exact rational
relations, with its conventions and limiting argument explained
in the corresponding section. A second certificate, using nine
of the same relation rows, evaluates the intermediate $h_{1,2}$
and reconciles the actual printed Family 2 right side exactly.

The JSON records specify the checks actually performed. Floating
point residuals are diagnostics, not certified error intervals.
An exact certificate verifies an algebraic reduction under its
stated analytic identities; it is not by itself a formalization
of convergence in a proof assistant. No decimal relation search
is promoted to a proof of $S_6$ or $S_8$.

\section{Further research questions}
\label{sec:research}

The results suggest several concrete next problems. The order below
prioritizes exact identities and structural classification over
estimates. Where a question is already in the supplied sources,
the proposed next step is stated relative to the new formulas.

\begin{enumerate}[label=\textbf{Q\arabic*.},leftmargin=*]

\item \textbf{Find the first non-Gaussian moment and its exact defect.}
For the statistic $X_{n,\sigma}$, determine
\[
 \mathbb E X_{n,\sigma}^{\,n+2}
       -\mathbb E Z^{\,n+2},\qquad Z\sim N(0,1),
\]
and decide whether $n+1$ is the exact universal matching degree.
The distinct-parameter proof stops because there are only
$n+1$ pair labels. Repeated powers of one label bring in
$S_3,S_5,\ldots$; the complete generating identity
\eqref{gn:eq:polygamma-germ} gives coupled equations for these
new coordinates. This is an algebraic obstruction to isolate,
not evidence that the next moment must be nonzero in every case.

\item \textbf{Classify joint polygamma moments.}
Give a finite combinatorial rule for the mixed cumulants of
$S_1,S_3,S_5,\ldots$ obtainable from the paired deformation.
The logarithm of the Gamma product involves interactions of at
most two parameter labels, but repeated powers of a label carry
higher polygamma statistics. A useful theorem would state
precisely which joint moments can be separated at dimension $n$,
and give their coefficients as polygamma values at $2\sigma$.
At $\sigma=1/4$, reflection should yield especially short
rational combinations of powers of $\pi$.

\item \textbf{Extend the determinant evaluation to later source rows.}
For the determinant in \eqref{gn:eq:F-explicit}, express a
prescribed family of later row sets as the evaluated Gustafson
normalizer times an explicit symmetric polynomial or a finite
hypergeometric sum. Such an identity would preserve the
oscillatory polynomial insertions exactly. It could address
the source's fixed-gap sign changes without an unjustified
total-positivity argument. The gap-six statement
$\det F_{b,3}>0$ for $b\ge7$ remains a concrete open target
\cite{repoordered}.

\item \textbf{Classify simultaneous zeros of colored tails.}
For
$A_k(\omega,a)=(\omega\partial_\omega+a)^k(1-\omega)^{-1}$,
classify all common zeros of consecutive $A_k$ at $a>0$ and
$\omega\in\Omega$. The example in this article proves that
consecutive zeros do occur. Determine the resulting dimension
of the space of homogeneous combinations with enhanced decay,
and the exact extra ordinary moments that each zero pattern
permits. A classification on the unit circle may be simpler
than one on the entire slit domain.

\item \textbf{Give a character-theoretic exceptional-moment dictionary.}
Apply finite combinations of root-of-unity colors to isolate
Dirichlet characters. At a negative integer where the relevant
$L$-function vanishes, determine the corresponding ordinary
Stieltjes-difference integral and simplify it using the
functional equation. The alternating beta family
\eqref{col:beta-formula} is the first case. A complete statement
should distinguish primitive and imprimitive characters and
retain their conductor and parity factors.

\item \textbf{Evaluate the first bilinear colored moment.}
For positive shifts, determine a natural ordinary domain and an
exact function class for
\[
 \int_0^\infty x^{u-1}
                  C_p^\omega(x,a)C_q^\eta(x,b)\dd x.
\]
The scalar Mellin formula suggests a Mellin convolution, but
the transform of a product is not the product of transforms.
One target is a two-parameter beta/Feynman integral whose
integer contacts reduce to a finite family of multiple-Lerch
coordinates. Coincident shifts and inverse colors are plausible
specializations to analyze first.

\item \textbf{Resolve the three-cluster spectral face.}
For $B^r,C,D$ with three distinct nonnegative masses, aggregate
the independent regulators into three sums before extracting
the squarefree coefficient. The two-cluster proof then suggests
a two-simplex hypergeometric integral and a finite iterated
integral dictionary. Determine the correct multivariable
polylogarithm arguments and branch domain. The answer should
explain what changes when $C=D$, where the depth-three formula
\eqref{tc:eq:two-displaced} must be recovered.

\item \textbf{Determine genuine depth reductions and lattice closures.}
The bound $\min(r,s)+1$ is a supplied representation bound.
Can the two-cluster combinations be reduced further as
functions? Independently, characterize parameters for which
the outer sum in \eqref{tc:eq:Gamma-dictionary} collapses to
finitely many multiple Gamma, Hurwitz, or polylogarithm
coordinates. These are two different problems: reducing
depth termwise does not eliminate an infinite spectral sum.

\item \textbf{Iterate anchored primitives with their constants intact.}
The colored primitive \eqref{col:primitive} lowers the Lerch
order by one; the cluster derivative deletes an initial index
one in the polylogarithm word. Iterate both operations and
identify a canonical Barnes-function normalization for the
root-of-unity cases. At every stage, retain explicit anchors.
Determine whether the normalized color-one limit commutes
with these primitives and identify any additional polynomial
terms when it does not.

\item \textbf{Supply a new relation for $S_6$ and revised $S_8$.}
The manuscript already distinguishes the conjectural period
equalities from its exact bounded-span and shuffle-obstruction
theorems. A proof must introduce a relation outside the
already excluded span, or a new analytic integral identity,
with all branch and endpoint terms included. The weight-three
worked reduction in this article is a useful small exact
certificate; it does not remove the higher-weight obstruction.
An appropriate deliverable is a rational certificate whose
primitive rows have independent analytic proofs, not only an
enlarged numerical relation search.

\item \textbf{Formalize the reusable analytic interfaces.}
The most useful first targets are continuous Cauchy--Binet with
dominated parameter differentiation, finite coefficient
extraction from normally convergent Mellin germs, and the
independent-regulator face identity
\eqref{tc:eq:independent-regulators}. Formalizing these
interfaces would allow many subsequent identities to be
checked by finite algebra while retaining the convergence
arguments that justify that algebra.

\end{enumerate}

